\documentclass[11pt]{article}

\usepackage[T1]{fontenc}
\usepackage[utf8]{inputenc}
\usepackage{amsmath,amssymb,amsthm}
\usepackage{booktabs}
\usepackage[margin=1in]{geometry}
\usepackage[hidelinks]{hyperref}
\usepackage{microtype}

\hypersetup{
  pdftitle={Exact Hurwitz Stability of Independence Polynomials for All Trees on 21 Vertices},
  pdfauthor={Xiaojian Zeng},
  pdfsubject={Exact computer-assisted tree enumeration and Routh--Hurwitz certification},
  pdfkeywords={independence polynomial, tree, Hurwitz stability, Routh test, exhaustive enumeration}
}

\newtheorem{theorem}{Theorem}
\newtheorem{lemma}[theorem]{Lemma}
\newtheorem{proposition}[theorem]{Proposition}

\title{Exact Hurwitz Stability of Independence Polynomials\\
for All Trees on 21 Vertices}
\author{Xiaojian Zeng\\
\small PBC School of Finance, Tsinghua University, Beijing 100190, China\\
\small \texttt{xiaojianzeng@foxmail.com}}
\date{September 7, 2026}

\begin{document}
\maketitle

\begin{abstract}
For a tree $T$, let $I(T,z)=\sum_k i_k(T)z^k$ be its independence
polynomial.  We give a computer-assisted exact proof that every zero of
$I(T,z)$ has negative real part for every tree on exactly 21 vertices.
The proof enumerates all $2{,}144{,}505$ unlabeled trees once, computes every
coefficient vector by integer tree dynamic programming, and applies an
integer-scaled Routh test.  Every Routh table is regular and has a strictly
positive first column; no singular or unstable case occurs.  A lossless
canonical-code census and a separately implemented systematic audit provide
checks on enumeration, polynomial construction, and stability certification.
Thus the closed-left-half-plane question at order 21 has an affirmative
answer, in the stronger strict sense.
\end{abstract}

\section{Statement and context}

For a finite simple graph $G$, an independent set is a vertex set containing
no edge, and
\[
  I(G,z)=\sum_{k=0}^{\alpha(G)}i_k(G)z^k
\]
counts independent sets by cardinality.  This polynomial was introduced in
this form by Gutman and Harary~\cite{GutmanHarary1983}.  Following Brown and
Cameron~\cite{BrownCameron2018}, we call it \emph{stable} when every zero lies
in the closed left half-plane.  We use \emph{strictly Hurwitz stable} when all
zeros lie in the open left half-plane.

Brown and Cameron proved stability for graphs of independence number at most
three and exhibited trees with roots arbitrarily far into the open right
half-plane.  They also reported a computation in which every tree of order at
most 20 was stable~\cite{BrownCameron2018}.  More recently, Zhang and Tu proved
that every spider is strictly Hurwitz stable~\cite{ZhangTu2026}.  Those results
do not decide the census at the next order.  The result here is limited to the
finite order-21 question and makes no general characterization or absolute
priority claim.

\begin{theorem}\label{thm:main}
For every finite simple tree $T$ on exactly 21 vertices, every zero of
$I(T,z)$ has strictly negative real part.  There are $2{,}144{,}505$ unlabeled
isomorphism classes in the exhaustive census, and none has a singular or
unstable Routh table.
\end{theorem}

The remainder gives the mathematical reductions and the exact finite
certificate underlying Theorem~\ref{thm:main}.

\section{Complete unlabeled-tree census}

The enumeration uses the free-tree generator of Wright, Richmond, Odlyzko,
and McKay~\cite{WrightRichmondOdlyzkoMcKay1986}.  Their algorithm extends the
Beyer--Hedetniemi rooted-tree successor and generates all nonisomorphic free
trees of a prescribed order without repetition.  Its objects are canonical
level sequences.  In our implementation, every vertex after the root is
joined to the nearest earlier vertex one level above it.  The resulting graph
is connected and has 20 edges, hence is a 21-vertex tree.

We added a generator-independent duplicate check.  Repeatedly remove leaves
to find the center or the two centers of a tree.  A rooted tree is encoded by
an opening bit, the recursively encoded child subtrees sorted by their
$(\text{order},\text{code})$ pairs, and a closing bit.  For a bicentered tree,
a virtual root is attached to the two center-rooted halves before the same
construction is applied.  Balanced-parenthesis parsing recursively recovers
the multiset of child encodings.  Consequently, equal codes are equivalent to
rooted isomorphism at each recursion depth, and the center convention makes
the final code injective on unrooted isomorphism classes.

At order 21 every code fits losslessly in a 64-bit unsigned integer.  Sorting
all emitted codes gave $2{,}144{,}505$ distinct values and no collision.  The
same total is listed independently in OEIS A000055~\cite{OEISA000055}; this
agreement is a cross-check, while enumeration completeness comes from the
generator theorem and the explicit implementation.

\section{Exact polynomial and Routh arithmetic}

Root a generated tree at the first vertex.  For each vertex $v$, let $A_v$
and $B_v$ denote the independence polynomials of the descendant subtree,
conditional on excluding and including $v$, respectively.

\begin{lemma}\label{lem:dp}
If the children of $v$ are $w$, then
\[
 A_v=\prod_w(A_w+B_w),\qquad
 B_v=z\prod_w A_w,
 \qquad I(T,z)=A_r+B_r.
\]
Thus the complete coefficient vector of $I(T,z)$ is obtained using integer
addition and multiplication only.
\end{lemma}

\begin{proof}
When $v$ is excluded, each child may independently be included or excluded.
When $v$ is included, every child is excluded, and the choice of $v$
contributes one factor of $z$.  Disjoint child subtrees make the contributions
multiplicative.  The root has no external condition, giving the last identity.
\end{proof}

As a universal low-degree check, every computed vector begins with
\[
  [i_0,i_1,i_2]=[1,21,\tbinom{21}{2}-20]=[1,21,190].
\]

For completeness, we describe the exact scaling used in the Routh tables.
Write coefficients in descending order and let $U$ and $L$ be two successive
Routh rows, with pivot $p=L_0$.  The ordinary next-row entry is
\begin{equation}\label{eq:routh}
  N_j=\frac{pU_{j+1}-U_0L_{j+1}}{p}.
\end{equation}
Instead of retaining rationals, the program stores
\begin{equation}\label{eq:integerrow}
  \operatorname{sgn}(p)
  \bigl(pU_{j+1}-U_0L_{j+1}\bigr)
\end{equation}
and divides the entire row only by its positive greatest common divisor.

\begin{lemma}\label{lem:scaling}
Until a zero pivot or an all-zero row is encountered, every stored row is a
positive scalar multiple of its ordinary rational Routh row.  In particular,
the stored first-column signs are the ordinary Routh signs.
\end{lemma}

\begin{proof}
The first two rows are unchanged.  Suppose the stored rows are positive
multiples $aU^*$ and $bL^*$ of the ordinary rows.  Then $p=bL^*_0$, and the
cross-product in~\eqref{eq:integerrow} is $ab$ times the corresponding
ordinary cross-product.  Multiplication by $\operatorname{sgn}(p)$ changes
division by $L^*_0$ into multiplication by the positive factor $|L^*_0|$.
The optional row-gcd division is also by a positive number.  Induction proves
the assertion.
\end{proof}

The Routh--Hurwitz criterion~\cite{Routh1877,Hurwitz1895} now gives an exact decision:
for a regular table the first-column sign variations count the roots in the
open right half-plane, while the singular cases that can signal an
imaginary-axis root are explicitly reported.  The constant coefficient of an
independence polynomial is 1, so zero itself cannot be a root.

\section{Exhaustive result and checks}

The exact run was divided into 43 committed chunks.  Each tree record entering
a chunk commitment contains its canonical code, full integer coefficient
vector, Routh status, and exact first-column signs.  The aggregate result is
shown in Table~\ref{tab:census}.

\begin{table}[ht]
\centering
\begin{tabular}{lr}
\toprule
Classification & Count\\
\midrule
generated unlabeled trees & $2{,}144{,}505$\\
distinct canonical codes & $2{,}144{,}505$\\
regular tables with positive first column & $2{,}144{,}505$\\
regular unstable tables & $0$\\
singular Routh tables & $0$\\
\bottomrule
\end{tabular}
\caption{Exact order-21 census.}\label{tab:census}
\end{table}

Lemma~\ref{lem:scaling} and Table~\ref{tab:census} show that every polynomial
has no right-half-plane root.  Because every table is regular, no
imaginary-axis case occurs.  This proves Theorem~\ref{thm:main}.

The primary run was checked in several different ways.
\begin{itemize}
  \item An integrity pass verified completion, contiguous chunk indices,
  absence of exception records, the expected binary length, and global
  uniqueness of all canonical codes.
  \item A separate NetworkX~3.4.2 graph-object stream regenerated the full
  count.  At the first and last indices and every multiple of 5,000---430
  predetermined cases---it matched canonical codes, recomputed $I(T,z)$ from
  the vertex-deletion recurrence
  \[
      I(G,z)=I(G-v,z)+zI(G-N[v],z),
  \]
  matched a separately rooted tree dynamic program, and found every leading
  Hurwitz determinant positive by fraction-free Bareiss elimination.
  \item A fresh complete execution reproduced all 43 chunk classifications
  and the entire canonical-code stream.  Small-order attacks also compared
  the generator with labeled Prüfer-sequence enumeration through order 7 and
  compared the polynomial dynamic program with brute-force subsets through
  order 12.
\end{itemize}

This separation between exhaustive production and reconstructing verification
is also used in recent exact finite graph work: for example, Zeng certifies
$6{,}021$ graph isomorphism classes with integer witnesses and a verifier that
reconstructs the relevant matrices~\cite{Zeng2026Whiskered}.  We cite that
paper as methodological context only; none of its mathematical conclusions is
used here.

\section*{Reproducibility and assistance disclosure}

The accompanying archive contains the decisive Python program, the 43 chunk
records, the complete canonical-code stream, the census integrity checker,
and the independently implemented NetworkX/Bareiss checker.  The principal
commands from the archive root are
\begin{verbatim}
python3 evidence/census21.py --order 21 --output evidence/census21
python3 evidence/audit_census21.py
python3 evidence/independent_checks.py
\end{verbatim}
The decisive polynomial and Routh computations use arbitrary-precision
integers; floating-point root approximations play no role in the theorem.

OpenAI Codex assisted with literature triage, implementation and manuscript
drafting, and a separate audit pass.  The theorem rests on the stated
mathematical reductions, exact machine-checkable computations, and reproduced
checks; no language-model assertion or confidence score is treated as a
mathematical certificate.

\bibliographystyle{abbrv}
\bibliography{references}

\end{document}
